\begin{abstract}
\noindent\textbf{Background:} Machine-learning models for type~2 diabetes mellitus (T2DM) screening often report optimistic performance because imputation or feature selection is fitted before validation, and because differences between models are rarely tested statistically.

\noindent\textbf{Methods:} We built a leakage-controlled pipeline that combines median imputation of impossible zero values, eight engineered clinical features, recursive feature elimination with cross-validation (RFECV) and an eight-learner stacking ensemble with a logistic-regression meta-learner. Every step was re-fitted inside each validation split. On the PIMA Indians Diabetes dataset (768 women, 268 with diabetes), we compared 11 learners with and without a common RFECV subset using 5$\times$10-fold repeated cross-validation with corrected confidence intervals and Holm-adjusted tests. Results were confirmed on a 154-patient hold-out set with bootstrap intervals and DeLong tests. We also assessed calibration, sensitivity-targeted thresholds, decision curves, alternative imputation strategies and SHAP explanations.

\noindent\textbf{Results:} RFECV reduced the inputs by 25\% without loss of discrimination; glucose, BMI, family history and their interactions were consistently retained. Cross-validated ROC-AUC was 0.843 (95\% CI 0.809--0.878) for RFECV + stacking and 0.842 (0.809--0.875) for logistic regression (adjusted $p = 1.00$). On the hold-out set, RFECV + stacking achieved sensitivity 0.815, specificity 0.740, MCC 0.532 and ROC-AUC 0.822. These were the highest observed sensitivity and MCC, but no ROC-AUC difference was significant. Class weighting caused systematic risk overestimation (calibration intercept $-0.62$, slope 1.00).

\noindent\textbf{Conclusions:} Under rigorous, leakage-free validation, stacking and RFECV give compact, sensitivity-oriented models but no discrimination gain over logistic regression on this small cohort. Calibration and threshold choice deserve as much attention as model choice.
\end{abstract}
